% ============================================================================
%  01-regression-from-scratch — writeup
%  Compile:  pdflatex writeup.tex   (run twice so \ref/\label resolve)
%  Or just paste into Overleaf. No special packages, no shell-escape needed.
% ============================================================================
\documentclass[11pt]{article}

% --- the only packages you need ---------------------------------------------
\usepackage[margin=1in]{geometry}   % sane margins
\usepackage{amsmath,amssymb}        % math
\usepackage{graphicx}               % \includegraphics for your figures
\usepackage{hyperref}               % clickable refs/links (optional; delete if you want)

% --- a few tiny convenience macros (delete any you don't use) ---------------
\newcommand{\R}{\mathbb{R}}
\newcommand{\norm}[1]{\left\lVert #1 \right\rVert}
\newcommand{\X}{\mathbf{X}}         % design matrix
\newcommand{\y}{\mathbf{y}}         % target vector
\DeclareMathOperator*{\argmin}{arg\,min}

\title{Econometrics Student Experiences the Machine Learning Framing of Regression, For Once in His Life}
\author{Kavi Mookherjee Amodt}
\date{\today}

\begin{document}
\maketitle

\begin{abstract}
Here I will do two things: (A) I will explain the machine learning framing of regression 
to someone who knows regression only as it is framed in econometrics (or in social sciences more generally), 
and (B), I will show how I implemented linear regression, logistic regression, and normalization 
entirely from scratch in Python. I write almost entirely in bullet points because that's how I like to learn. Each bullet deepens understanding; all of them can be seen as having approximately equal semantic weight, so that you don't have to decide whether what you are reading is fluff or not; and you can pace yourself metronomically, it being easy to return to your last checkpoint of understanding. 
\end{abstract}

\vspace{0.5cm}

% ============================================================================
{\huge Part (A): Regression in ML and in the Social Sciences}

\section{Inference vs prediction in statistics}\label{sec:cultures}
    \subsection{In econometrics we infer}
        \begin{itemize}
            \item In econometrics, regression is the engine of \textbf{inference}.
            \item Inference, in this sense, means we use existing data to infer the \textit{relationship} between some variables (e.g. $x$ = income and $y$ = SAT score).
            \item To be more precise, this relationship is a statistical distribution. Our goal is to model that distribution. 
            \item When we do regression, we assume that there's a \textit{true} distribution; a \textbf{data generating process}.
            \item When we do linear regression, we assume that the true data generating process looks like a (multivariate) linear equation: \[y = X\beta + \epsilon\]
            \item Our \textbf{object of interest} is the vector of weights: $\beta$.
            \item Put differently, $\beta$ is the vector of true relationships between our variables. $\epsilon$ is the noise (more on this later).
            \item Our goal is to create an estimate of the coefficient vector $\beta$—that estimate is denoted $\hat{\beta}$—which comes closest, when multiplied by the \textbf{design matrix} $X$, to the data for $y$: \[\hat{y} = X\hat{\beta}\]
        \end{itemize} 

    \subsection{In ML we predict}
        \begin{itemize}
            \item In machine learning, our primary goal is not to obtain and report a relationship (like $\hat{\beta}$), but rather to use data from the world to \textit{make more data}.
            \item In other words, we take a complex set of causal relationships out there in the world—a set of samples from a data generating process—and use it to \textbf{predict} what that process would output.
            \item An LLM doesn't explicitly mimick the causal process in your brain which leads to language (emotional disregulation, excitement, boredom, etc); it finds complex statistical \textit{relationships} between tokens (samples from the process) and predicts what the causal process in your brain would output (the next word).
            \item Thus, our object of interest is not $\beta$, but $y$. 
            \item Instead of focusing on obtaining and reporting $\hat{\beta}$, we try to make a prediction: $\hat{y}$.
        \end{itemize}

    \subsection{Throw $\epsilon$ out the window (for now)}
        \begin{itemize}
            \item FINISH THIS
        \end{itemize}

\section{Linear regression is just projecting onto span(X)}\label{sec:projection}
    \subsection{Gradient of the squared-error loss function}
        \begin{itemize}
            \item Notice how, in the ML framing, there is no notion of probability. It's just pure optimization of this loss function: \[L(\beta) = \frac{1}{n}||y - X\beta||^2\]
            \item Remember that $||v|| = v^Tv$ and $(A - B)^T = A^T - B^T$
            \item Therefore, in matrix form: $$L(\beta) \quad = \quad (y - X\beta)^T(y - X\beta) \quad = \quad y^Ty - 2\beta^TX^Ty + \beta^TX^TX\beta$$
            \item Thus, setting $\nabla L = 0 \implies -2X^Ty + 2X^TX\beta = 0 \implies \boldsymbol{\hat{\beta}} \; \mathbf{ = (X^TX)^{-1}X^Ty}$    
            \item This result is called the \textbf{normal} equation; we use this to compute our estimate of the relationship vector $\hat{\beta}$.
            \item Crucially, this is the \textit{closed-form} way to compute $\hat{\beta}$ for linear regression—we don't have such a luxury in logisitic regression.          
        \end{itemize}

    \subsection{H, the Hat matrix}
        \begin{itemize}
            \item Plugging our normal equation for $\hat{\beta}$ back into the equation for our $\hat{y}$ prediction: \[\hat{y} = X\beta \quad = \quad X(X^TX)^{-1}X^Ty \quad = \quad \mathbf{H}y\]
            \item $\mathbf{H}$ is the Hat matrix of $X$: multiplying a vector $y$ by the Hat matrix of $X$ gives you the \textbf{orthogonal projection} of $y$ onto $X$.
            \item That's all OLS is: our estimator for $y$ is the point in span($X$) which is closest, in terms of Euclidean distance, to $y$. 
            \item That mechanic right there—finding the closest point in representable space to our datapoint—is one you will see throughout linear-algebraic machine learning: in QR decomposition, principal components analysis, support vector machines.
        \end{itemize}

\section{Gradient descent for OLS}\label{sec:GD}
    \subsection{Why use gradient descent?}
        \begin{itemize}
            \item If we have a closed-form way to compute our coefficients, why is gradient descent considered the \textit{sledgehammer} of modern ML?
            \item Unsurprisingly, the answer is efficiency: computing $X^TX$ is $O(p^3)$ where $p$ is the number of predictors.
            \item We thus use GD in both ML and econometric settings.
        \end{itemize}

    \subsection{How GD works}
        \begin{itemize}
            \item The gradient of a continuous, multivariate function (in this case, our loss function $L$) is a vector which points in the direction of \textbf{steepest ascent}.
            \item As we saw earlier, it is denoted $\nabla L$. Its magnitude $||\nabla L||$ equals the slope of L in that direction. 
            \item GD is a process:
                \begin{enumerate}
                    \item Find $\nabla L$ at the point on the ``surface" of the function L where you stand now. 
                    \item The vector $\hat{\beta} \leftarrow \hat{\beta} - \mu\nabla L$.
                    \item Travel to the t
                \end{enumerate}
        \end{itemize}

\section{OLS is MLE under Gaussian noise}\label{sec:MLE}

\section{Logistic regression doesn't have closed form}\label{sec:logistic}

\section{Regularizing is biasing on purpose}\label{sec:regularizing}

\vspace{0.5cm}

{\huge Part (B): Implementation Details}



\end{document}